\chapter{Speech development in children}
\index{Speech development in children}
\medskip
\noindent\textit{Educational draft; seek professional advice for individual care.}

\section*{Definition and scope}
Speech development in children concerns the health, development, or safety of infants, children, or adolescents. Assessment must account for age, development, communication, family context, and safeguarding.

\section*{Contributors and risk}
Consider infection, congenital or genetic factors, developmental variation, injury, exposure, nutrition, and emotional or social stress.

\section*{Presentation}
Record onset, progression, fever or pain, feeding, hydration, breathing, sleep, behaviour, development, school function, injury, and exposures.

\section*{Assessment}
Use examination, growth and developmental review, hearing or vision assessment, and targeted tests appropriate to age.

\section*{Management}
Management may involve observation, supportive care, therapy, medicines, procedures, or paediatric and allied-health referral.

\section*{When to seek urgent care}
Seek urgent care for breathing difficulty, blue or very pale appearance, seizure, reduced consciousness, severe dehydration, non-blanching rash, serious injury, or safeguarding concern.

\section*{Prevention and follow-up}
Use immunisation, safe sleep and transport practices, supervision, nutrition, and routine health checks.
\section*{Safety-netting}
Seek urgent help for severe or rapidly worsening symptoms, collapse, breathing difficulty, severe pain, confusion, dehydration, uncontrolled bleeding, or any immediate danger.
section*{Sources and limitations}
This concise educational overview should be checked against current local clinical guidance and adapted to the individual.
